\documentclass[12pt,a4paper]{article}

\usepackage[margin=2.2cm]{geometry}
\usepackage[ngerman]{babel}
\usepackage[T1]{fontenc}
\usepackage[utf8]{inputenc}
\usepackage{lmodern}
\usepackage{array}
\usepackage{longtable}
\usepackage{booktabs}
\usepackage[table]{xcolor}
\usepackage{enumitem}
\usepackage{amsmath}
\usepackage{fancyhdr}
\usepackage{hyperref}

\pagestyle{fancy}
\fancyhf{}
\lhead{Physik -- Einführung in den Drehimpuls}
\rhead{Lektionsablauf}
\cfoot{\thepage}

\setlength{\parindent}{0pt}
\setlength{\parskip}{6pt}

\definecolor{headgray}{RGB}{230,230,230}
\definecolor{lightrow}{RGB}{245,245,245}

\begin{document}

\begin{center}
    {\LARGE \textbf{Lektionsablauf: Einführung in den Drehimpuls}}\\[0.3cm]
    \large 6. Maturitätsklasse -- 45 Minuten
\end{center}

\vspace{0.4cm}

\textbf{Thema:} Einführung in den Drehimpuls aufbauend auf dem Vorwissen zu Kreisbewegung, Winkelgeschwindigkeit, Drehmoment und translatorischem Impuls.

\textbf{Lernziele:}
\begin{itemize}[leftmargin=1.2cm]
    \item Die Schülerinnen und Schüler knüpfen an das Vorwissen zu Impuls, Kreisbewegung und Drehmoment an.
    \item Sie beschreiben den Drehimpuls als Rotationsanalogon des translatorischen Impulses.
    \item Sie deuten das Drehstuhl-Experiment mithilfe der Erhaltung des Drehimpulses.
    \item Sie wenden die Beziehung \(L = I\omega\) qualitativ auf einfache Situationen an.
\end{itemize}

\textbf{Material:}
\begin{itemize}[leftmargin=1.2cm]
    \item Drehstuhl
    \item zwei Hanteln oder Gewichte
    \item Tafel / Beamer
    \item Arbeitsblatt
    \item Clicker-System / Abstimmungstool
\end{itemize}

\vspace{0.3cm}

\renewcommand{\arraystretch}{1.4}
\setlength{\LTpre}{6pt}
\setlength{\LTpost}{6pt}

\begin{longtable}{@{}p{1.8cm}p{5.9cm}p{4.2cm}p{2.6cm}p{2.5cm}@{}}
\rowcolor{headgray}
\textbf{Zeit} & \textbf{Aktivität / Inhalt} & \textbf{Lehrperson} & \textbf{Schülerinnen und Schüler} & \textbf{Sozialform} \\
\midrule
\endfirsthead

\rowcolor{headgray}
\textbf{Zeit} & \textbf{Aktivität / Inhalt} & \textbf{Lehrperson} & \textbf{Schülerinnen und Schüler} & \textbf{Sozialform} \\
\midrule
\endhead

\bottomrule
\endfoot

\rowcolor{lightrow}
0--3 min.
&
\textbf{Einstieg und Fokussierung} \newline
Begrüßung, Thema der Lektion und Leitfrage: \glqq Warum wird eine Drehbewegung manchmal schneller, obwohl niemand zusätzlich anschiebt?\grqq
&
Begrüßt die Klasse, stellt das Thema vor und aktiviert Aufmerksamkeit mit der Leitfrage.
&
Hören zu, nehmen die Leitfrage auf, richten Material und Arbeitsblatt bereit.
&
Plenum
\\

3--8 min.
&
\textbf{Demonstrationsexperiment: Drehstuhl} \newline
Eine Schülerin oder ein Schüler sitzt auf dem Drehstuhl mit ausgestreckten Armen und Gewichten in den Händen. Die Person wird langsam in Rotation versetzt und zieht dann die Arme ein.
&
Instruiert das Experiment, sorgt für Sicherheit, führt oder begleitet die Demonstration und stellt Beobachtungsfragen: \glqq Was verändert sich? Was nicht?\grqq
&
Beobachten das Experiment, notieren ihre Beobachtungen auf dem Arbeitsblatt und formulieren erste Vermutungen.
&
Plenum / Beobachtung
\\

8--12 min.
&
\textbf{Sammeln von Hypothesen} \newline
Kurze Auswertung des Experiments: mögliche Erklärungen und spontane Ideen werden gesammelt.
&
Moderiert das Sammeln an der Tafel, greift typische Vorstellungen auf, ohne sie sofort zu bewerten.
&
Nennen Vermutungen, z.\,B. kleinere Reibung, Gewichte ziehen nach innen, schnellere Bewegung wegen kleinerem Radius.
&
Plenum
\\

\rowcolor{lightrow}
12--17 min.
&
\textbf{Aktivierung des Vorwissens} \newline
Wiederholung zentraler Begriffe:
\begin{itemize}[nosep,leftmargin=0.5cm]
    \item Kreisbewegung
    \item Winkelgeschwindigkeit \(\omega\)
    \item translatorischer Impuls \(p = mv\)
    \item Drehmoment \(\tau\)
\end{itemize}
&
Stellt gezielte Fragen, lässt bekannte Beziehungen nennen und schreibt die Schlüsselbegriffe strukturiert an die Tafel.
&
Beantworten Fragen, ergänzen Definitionen im Arbeitsblatt und verknüpfen bekanntes Wissen mit dem Experiment.
&
Plenum
\\

17--21 min.
&
\textbf{Überleitung zum neuen Begriff} \newline
Vom translatorischen Impuls zur Rotationsbewegung: \glqq Wenn der Impuls die Bewegungsgröße in der Translation ist, wie heißt die passende Größe bei Rotationen?\grqq
&
Führt die Analogie Translation -- Rotation ein und motiviert den Begriff \textbf{Drehimpuls}.
&
Denken mit, bringen Vorschläge ein und diskutieren kurz, welche Größe bei Rotationen analog zum Impuls sein könnte.
&
Plenum / Unterrichtsgespräch
\\

21--26 min.
&
\textbf{Erarbeitung des Drehimpulses} \newline
Einführung der Beziehungen
\[
\vec L = \vec r \times \vec p,
\qquad
L = mvr,
\qquad
L = I\omega
\]
für den Spezialfall der Kreisbewegung.
&
Erklärt die Bedeutung der Formeln schrittweise, stellt die Verbindung zu \(v=\omega r\) her und betont die physikalische Idee statt formaler Herleitung im Detail.
&
Füllen das Arbeitsblatt aus, notieren Begriffe und stellen Verständnisfragen.
&
Plenum mit gelenkter Einzelarbeit
\\

\rowcolor{lightrow}
26--31 min.
&
\textbf{Clicker-Aktivität} \newline
Konzeptfrage, z.\,B.: \glqq Eine Person dreht sich auf einem Drehstuhl und zieht die Arme ein. Was passiert mit Drehimpuls, Trägheitsmoment und Winkelgeschwindigkeit?\grqq
&
Projiziert die Frage, gibt Antwortoptionen vor, startet die Abstimmung und lässt kurz begründen.
&
Denken zunächst individuell nach, stimmen mit dem Clicker ab und vergleichen danach ihre Überlegungen mit der Klasse.
&
Einzelarbeit $\rightarrow$ Plenum
\\

31--34 min.
&
\textbf{Auswertung der Clicker-Frage} \newline
Besprechung der Abstimmung, Klärung typischer Fehlvorstellungen.
&
Greift falsche und richtige Antworten auf, erklärt präzise: ohne äußeres Drehmoment bleibt \(L\) konstant; beim Einziehen der Arme wird \(I\) kleiner und \(\omega\) größer.
&
Begründen ihre Wahl, korrigieren gegebenenfalls ihre Vorstellung und ergänzen die Lösung auf dem Arbeitsblatt.
&
Plenum
\\

34--39 min.
&
\textbf{Rückbezug auf das Experiment} \newline
Das Experiment wird nun mit dem neu eingeführten Konzept erklärt:
\[
L = I\omega = \text{konstant}
\]
&
Verknüpft die Theorie mit der beobachteten Demonstration und formuliert den Erhaltungsgedanken explizit.
&
Erklären in eigenen Worten, weshalb sich die Person nach dem Einziehen der Arme schneller dreht.
&
Plenum / kurzes Unterrichtsgespräch
\\

\rowcolor{lightrow}
39--43 min.
&
\textbf{Sicherung} \newline
Gemeinsame Zusammenfassung der Kernaussagen:
\begin{itemize}[nosep,leftmargin=0.5cm]
    \item Drehimpuls als Rotationsanalogon zum Impuls
    \item \(L = I\omega\)
    \item äußeres Drehmoment verändert \(L\)
    \item ohne äußeres Drehmoment bleibt \(L\) erhalten
\end{itemize}
&
Leitet die Sicherung, strukturiert die Merksätze an der Tafel und fordert kurze Schülerformulierungen ein.
&
Formulieren zentrale Erkenntnisse, ergänzen Merksätze im Arbeitsblatt.
&
Plenum
\\

43--45 min.
&
\textbf{Abschluss / Ausblick} \newline
Kurzer Transfer: Eiskunstlauf, Kreisel oder Astronomie als Anwendungen des Drehimpulses.
&
Gibt einen Ausblick auf weitere Anwendungen und schließt die Lektion.
&
Hören zu, stellen letzte Fragen und geben das Arbeitsblatt ab oder behalten es als Zusammenfassung.
&
Plenum
\\

\end{longtable}

\vspace{0.3cm}

\section*{Didaktische Begründung in Kürze}

Der Ablauf beginnt mit einem \textbf{anschaulichen Experiment}, das ein kognitives Problem erzeugt und damit die Motivation für den neuen Begriff schafft. Anschließend wird das \textbf{Vorwissen gezielt aktiviert}, damit die Schülerinnen und Schüler den Drehimpuls nicht als isolierte Formel, sondern als \textbf{Rotationsanalogon des bekannten Impulsbegriffs} verstehen. Die \textbf{Clicker-Aktivität} dient der Konzeptüberprüfung und macht Fehlvorstellungen sichtbar. Zum Schluss wird das Experiment mithilfe der Erhaltung des Drehimpulses erklärt und die zentrale Idee gesichert.

\section*{Vorschlag für die Clicker-Frage}

\textbf{Frage:} Eine Person sitzt auf einem Drehstuhl und dreht sich mit ausgestreckten Armen. Danach zieht sie die Arme ein. Welche Aussage ist richtig?

\begin{enumerate}[label=\Alph*.]
    \item Der Drehimpuls wird größer, weil sich die Person schneller dreht.
    \item Das Trägheitsmoment wird kleiner, die Winkelgeschwindigkeit größer und der Drehimpuls bleibt gleich.
    \item Das Drehmoment wird größer, deshalb dreht sich die Person schneller.
    \item Die Masse wird kleiner, deshalb nimmt die Winkelgeschwindigkeit zu.
\end{enumerate}

\textbf{Richtige Antwort:} B

\section*{Hinweise zum Experiment}

\begin{itemize}[leftmargin=1.2cm]
    \item Sicherheit vorführen: stabile Sitzposition, langsame Rotation, Gewichte kontrolliert halten.
    \item Die Klasse vor dem Einziehen der Arme ausdrücklich zur Vorhersage auffordern.
    \item Beobachtung und Erklärung klar trennen: erst beschreiben lassen, dann deuten.
\end{itemize}

\end{document}